\section*{Module Overview}

\begin{keybox}[Introduction]
This module introduces the principles and practices of Human-AI Interaction (HAI), building on traditional Human-Computer Interaction (HCI) foundations while addressing the unique challenges posed by AI systems. The focus is on \textbf{applying} concepts to real design problems---not just memorizing definitions, but being able to analyze systems, identify issues, and propose solutions.
\end{keybox}

\begin{keybox}[Module Aims]
By the end of this module, you will be able to:
\begin{itemize}
  \item Understand and apply core interaction design concepts (affordances, signifiers, mental models).
  \item Analyze systems using HCI frameworks and heuristics.
  \item Design and evaluate user-centered AI systems.
  \item Understand the unique challenges of AI systems: automation, fairness, transparency, harms.
  \item Apply prototyping and evaluation methods appropriate to different design contexts.
  \item Critically assess AI systems for usability, fairness, and environmental impact.
\end{itemize}
\end{keybox}

\begin{keybox}[Exam Approach]
The exam tests your ability to \textbf{apply} concepts, not just recall definitions:
\begin{itemize}
  \item You will be given a \textbf{context} (e.g., ``Design a vaccination passport app for NHS'').
  \item You must show you can use the concepts to analyze and solve problems in that context.
  \item Example tasks: Create interview questions for user needs; identify gulfs of evaluation/execution on a screenshot; suggest prototype approaches; outline evaluation metrics.
  \item There is no single ``right answer''---reasonable, well-justified responses that show understanding are valued.
\end{itemize}
\end{keybox}

\subsection*{Course Structure}

\begin{center}
\renewcommand{\arraystretch}{1.3}
\begin{tabular}{@{} c l p{7cm} @{}}
\toprule
\textbf{Week} & \textbf{Topic} & \textbf{Key Concepts} \\
\midrule
1 & Introduction to Interaction Design & Affordances, signifiers, interaction design goals, AI as design material \\
2 & Cognition and Mental Models & Perception, mental models (conceptual/represented/mental), gulfs of evaluation/execution, distributed cognition, CASA theory \\
3 & Usability Goals and Heuristics & Usability goals, Nielsen's heuristics, human-AI guidelines \\
4 & Human-Centered Design & Design lifecycles (Double Diamond, IDEO, Stanford d.school), needs vs.\ requirements, interview methods, empathy maps, user-centered AI development \\
5 & Prototyping & Prototype definition/function, assumption testing, fidelity levels, pros/cons \\
6 & \textit{Reading Week} & --- \\
7 & Automation, Agency, and Fairness & Control, automation levels, fairness definitions, AI harms, transparency vs.\ explainability \\
8 & Environmental Sustainability & AI energy impacts, sustainability calculations, sources of environmental impact \\
9 & Evaluation & Evaluation purpose, method selection, heuristic evaluation, think aloud, experiments (IV/DV, reliability, validity, between/within subjects) \\
10 & Red Teaming and Field Studies & Appropriation, technology probes, red teaming, attack types, defenses (constitutional classifiers) \\
\bottomrule
\end{tabular}
\end{center}

\subsection*{Key Lists to Remember}

The instructor emphasizes three key lists worth memorizing:
\begin{enumerate}
  \item \textbf{Usability Goals} (Week 3): ~6 goals
  \item \textbf{Nielsen's Heuristics} (Week 3): 10 heuristics
  \item \textbf{Types of AI Harms} (Week 7): Categories of potential harms from AI systems
\end{enumerate}

\subsection*{Core Exam Skills}

\begin{itemize}
  \item \textbf{Apply concepts to given contexts:} Analyze screenshots, design briefs, system descriptions.
  \item \textbf{Connect concepts across weeks:} Mental models $\leftrightarrow$ affordances $\leftrightarrow$ control $\leftrightarrow$ usability.
  \item \textbf{Create plausible designs:} Interview questions, prototypes, evaluation plans.
  \item \textbf{Argue your choices:} Justify why you chose a particular method or approach.
\end{itemize}

